% Источники: raw/IMG_7636.jpeg -- raw/IMG_7645.jpeg.
% Дубликаты: 7638 повторяет 7636, 7639 повторяет 7637, 7644 повторяет 7643.
% Компиляция: tectonic 08.tex
\ifdefined\lessonpreambleloaded
  \expandafter\endinput
\fi
\let\lessonpreambleloaded\empty
\documentclass[a4paper,11pt]{article}
\usepackage{fontspec}
\setmainfont{Times New Roman}
\usepackage[russian]{babel}
\usepackage{amsmath,amssymb}
\usepackage[margin=22mm]{geometry}
\usepackage{indentfirst}
\usepackage{needspace}
\usepackage{tikz}
\usetikzlibrary{patterns}
\usepackage[hidelinks]{hyperref}
\usepackage{fancyhdr}
\pagestyle{fancy}
\fancyhf{}
\renewcommand{\headrulewidth}{0pt}
\fancyfoot[C]{\footnotesize\url{https://github.com/bmstu-software-engineering/cs-s5-tv-26y}}
\fancyfoot[R]{\thepage}
\setlength{\emergencystretch}{2em}
\newcommand{\sourcepage}[1]{\par\Needspace{6\baselineskip}\medskip\noindent{\footnotesize\itshape Источник: \detokenize{#1}.}\par\nopagebreak\smallskip}
\newcommand{\editorial}[1]{\par\smallskip\noindent{\small\itshape Примечание к переносу. #1}\par\smallskip}
\newcommand{\lessontitle}[3]{\begin{center}{\large\bfseries Теория вероятностей. #1}\\[4pt]{\large #2}\\[4pt]#3\end{center}}
\newcommand{\parttitle}[1]{\par\Needspace{4\baselineskip}\medskip\noindent\textbf{#1}\par\nopagebreak\smallskip}
\tikzset{element/.style={draw,rectangle,minimum width=7mm,minimum height=5mm,fill=white},every picture/.style={line width=.5pt}}

\begin{document}
\lessontitle{Семинар №\,8}{Условная вероятность. Формула полной вероятности}{24.10.26}

\sourcepage{IMG_7636.jpeg; повторный снимок: IMG_7638.jpeg}
\section*{Условная вероятность}
\[
 P(A\mid B)=\frac{P(A\cap B)}{P(B)},\qquad P(B)>0.
\]

\section*{Три игральные кости}
Наудачу бросают три кости. Событие $A$ --- на всех костях выпали разные числа;
$B$ --- хотя бы на одной кости выпало $6$. Найти $P(A\mid B)$.

\parttitle{Способ 1. Через вероятности событий}
Исход --- упорядоченная тройка чисел от $1$ до $6$:
\[
 \Omega=\{(x_1,x_2,x_3):x_i\in\{1,\ldots,6\}\},\qquad N=6^3.
\]
Событию $\overline B$ соответствуют исходы, в которых не выпало ни одной шестёрки:
\[
 \overline B=\{(x_1,x_2,x_3):x_i\in\{1,\ldots,5\}\},\qquad
 N_{\overline B}=5^3.
\]
\[
 P(B)=1-P(\overline B)=1-\frac{5^3}{6^3}.
\]

\sourcepage{IMG_7637.jpeg; повторный снимок: IMG_7639.jpeg}
Для $A\cap B$ шестёрка выпадает ровно на одной кости, а два других числа
выбираются из $\{1,\ldots,5\}$ и различны:
\begin{align*}
 A\cap B={}&\{(6,u,v):u,v\in\{1,\ldots,5\},\ u\ne v\}\\
 &\cup\{(u,6,v):u,v\in\{1,\ldots,5\},\ u\ne v\}\\
 &\cup\{(u,v,6):u,v\in\{1,\ldots,5\},\ u\ne v\}.
\end{align*}
\[
 N_{A\cap B}=3A_5^2=3\frac{5!}{3!}=60,\qquad
 P(A\cap B)=\frac{60}{6^3}.
\]
\[
 P(A\mid B)=\frac{P(A\cap B)}{P(B)}.
\]
\editorial{На этой странице в подстановке потерян знаменатель $P(B)$ и получено
$5/18$. На следующей странице автор пересчитывает $B$ и исправляет ответ на $60/91$.}

\parttitle{Способ 2. Сужение пространства исходов}
При условии $B$ новым пространством становится $\Omega_B=B$.
В нём благоприятны $N_{A\cap B}=60$ исходов:
\[
 P(A\mid B)=\frac{N_{A\cap B}}{N_B}.
\]
\begin{center}
\begin{tikzpicture}[scale=.8]
 \draw (-1.8,-1.4) rectangle (3,1.4);
 \fill[gray!20] (1.2,0) circle (1);
 \begin{scope}\clip (0,0) circle (1);\fill[gray!50] (1.2,0) circle (1);\end{scope}
 \draw (0,0) circle (1) (1.2,0) circle (1);
 \node at (-.5,0){$A$};\node at (1.7,0){$B$};\node at (1.2,-1.15){$B=\Omega_B$};
 \node at (-1.45,-1.1){$\Omega$};
\end{tikzpicture}
\end{center}

\sourcepage{IMG_7640.jpeg}
Перечислим случаи для события $B$:
\begin{align*}
 B={}&\{(6,u,v):u,v\in\{1,\ldots,5\}\}
 \cup\{(u,6,v):u,v\in\{1,\ldots,5\}\}\\
 &\cup\{(u,v,6):u,v\in\{1,\ldots,5\}\}\\
 &\cup\{(6,6,u),(6,u,6),(u,6,6):u\in\{1,\ldots,5\}\}
 \cup\{(6,6,6)\}.
\end{align*}
\[
 N_B=25+25+25+5\cdot3+1=91,\qquad
 \boxed{P(A\mid B)=\frac{60}{91}.}
\]
Для обратной условной вероятности в исходнике записано:
\[
 P(B\mid A)=\frac{P(A\cap B)}{P(A)},\qquad
 P(A)=\frac{A_6^3}{6^3}=\frac{6!}{3!6^3}.
\]

\section*{Зависимые и независимые события}
Если $P(B)>0$, то независимость $A$ и $B$ равносильна
\[
 P(A\mid B)=P(A).
\]
При $P(A)>0$ эквивалентная запись:
\[
 P(B\mid A)=P(B).
\]
События $A_1,\ldots,A_n$ попарно независимы, если независимы любые
два события $A_i$ и $A_j$ при $i\ne j$.

\sourcepage{IMG_7641.jpeg}
\section*{Карта из колоды}
Из колоды в $52$ карты извлекают одну карту. В исходнике обозначены события:
\begin{itemize}
 \item $A$ --- извлечённая карта туз;
 \item $B$ --- карта чёрной масти;
 \item $C$ --- карта-картинка.
\end{itemize}
Требуется проверить попарную независимость. На странице записаны только
числа исходов и отдельные вероятности:
\[
 N=52,\qquad N_A=4,\qquad N_B=26,\qquad N_C=16,
\]
\[
 P(A)=\frac4{52},\qquad P(B)=\frac12,\qquad P(C)=\frac{16}{52}=\frac4{13}.
\]
\editorial{Число $16$ сохранено из исходника. Состав события «карта-картинка»
не уточняется: если это только валет, дама и король, таких карт $12$.
Проверка независимости на фотографии не закончена.}

\sourcepage{IMG_7642.jpeg}
\section*{Иностранные языки}
Из $100$ студентов английский знают $50$, французский --- $40$,
немецкий --- $35$; английский и французский --- $20$,
французский и немецкий --- $8$.
Нужно исследовать пары событий на зависимость и найти $P(A\cap B\mid C)$.
Наудачу выбирают одного студента:
\[
 A=\{\text{знает английский}\},\quad
 B=\{\text{знает французский}\},\quad C=\{\text{знает немецкий}\}.
\]
\[
 P(A)=\frac12,\quad P(B)=\frac{40}{100}=\frac25,\quad
 P(C)=\frac{35}{100}.
\]
На странице начаты вычисления:
\[
 P(A\mid B)=\frac{P(A\cap B)}{P(B)}
 =\frac{20/100}{40/100}=\frac12,
\]
\[
 P(A\mid C)=\frac{P(A\cap C)}{P(C)},\qquad
 P(B\mid C)=\frac{P(B\cap C)}{P(C)}.
\]
\editorial{Эти вычисления и дополнительный вопрос в исходнике не завершены.
Число знающих английский и немецкий, а также число знающих все три языка
на фотографии не указаны. Недостающие данные не восстанавливаются.}

\sourcepage{IMG_7643.jpeg; повторный снимок: IMG_7644.jpeg}
\section*{Формула полной вероятности}
Гипотезы $H_1,\ldots,H_n$ попарно несовместны и образуют полную группу:
\[
 H_i\cap H_j=\varnothing\quad(i\ne j),\qquad
 \bigcup_{i=1}^n H_i=\Omega.
\]
Для гипотез с положительными вероятностями
\[
 \boxed{P(A)=\sum_{i=1}^n P(H_i)P(A\mid H_i).}
\]

\section*{Телевизоры трёх заводов}
Доли телевизоров трёх заводов составляют $30\%$, $20\%$, $50\%$.
Доля брака на первом заводе --- $20\%$, на втором --- $10\%$,
на третьем --- $5\%$. Найти вероятность приобретения исправного телевизора.

$A$ --- приобретён исправный телевизор;
$H_i$ --- телевизор с $i$-го завода. Гипотезы попарно несовместны,
$H_1\cup H_2\cup H_3=\Omega$:
\[
 P(H_1)=\frac3{10},\qquad P(H_2)=\frac2{10},\qquad P(H_3)=\frac5{10},
 \qquad \sum_{i=1}^3P(H_i)=1.
\]

\sourcepage{IMG_7645.jpeg}
\[
 P(A\mid H_1)=0{,}8,\qquad P(A\mid H_2)=0{,}9,\qquad
 P(A\mid H_3)=0{,}95.
\]
\begin{align*}
 P(A)&=P(H_1)P(A\mid H_1)+P(H_2)P(A\mid H_2)+P(H_3)P(A\mid H_3)\\
 &=0{,}3\cdot0{,}8+0{,}2\cdot0{,}9+0{,}5\cdot0{,}95=0{,}895.
\end{align*}

\section*{Порядок вытягивания экзаменационного билета}
В записи рассматриваются $30$ билетов, из них $6$ --- «счастливые».
Исследуется, влияет ли порядок вытягивания билета на вероятность успеха Иванова.
$A$ --- Иванов вытащил счастливый билет.
\[
 P(A)=\frac6{30}=\frac15.
\]
Если до него билет уже вытянул другой человек, обозначим:
$H_1$ --- тот вытащил счастливый билет, $H_2$ --- несчастливый.
\[
 P(H_1)=\frac15,\qquad P(H_2)=\frac45.
\]
В первом случае остаются $5$ счастливых билетов из $29$,
во втором --- $6$ из $29$. Поэтому
\[
 P(H_1)P(A\mid H_1)=\frac15\frac5{29}=\frac1{29},\qquad
 P(H_2)P(A\mid H_2)=\frac45\frac6{29},
\]
\[
 P(A)=\frac15\frac5{29}+\frac45\frac6{29}=\frac15.
\]
Порядок вытягивания не влияет на вероятность счастливого билета.
\editorial{В исходнике произведения $P(H_i)P(A\mid H_i)$ подписаны как
условные вероятности. Обозначения уточнены; сами множители сохранены.}
\end{document}
